\documentclass[twoside]{article}
\setlength{\oddsidemargin}{0.2 in}
\setlength{\evensidemargin}{0.2in}
\setlength{\topmargin}{-0.8 in}
\setlength{\textwidth}{6.5 in}
\setlength{\textheight}{8.5 in}
\setlength{\headsep}{0.75 in}
\setlength{\parindent}{0 in}
\setlength{\parskip}{0.1 in}

%\usepackage{amsmath}
%\usepackage{amsthm}
\usepackage{amssymb}
\usepackage{amsthm}
\usepackage[tbtags]{amsmath}
\newtheorem{thm}{Theorem}
\usepackage[english]{babel}
\usepackage[utf8]{inputenc}
\usepackage{graphicx}
\usepackage[colorinlistoftodos]{todonotes}
\usepackage{hyperref}
\usepackage{setspace}

\begin{document}
\begin{center}
\begin{spacing}{1.75}
{\Large { ECE368: Probabilistic Reasoning} \\
	{Lab 2: Classification with Gaussian Models and Bayesian Linear Regression}}
\end{spacing}
\vspace{-1em}
\end{center}
%\vspace{-11mm}
You can complete this lab in a group of two. Please provide the name and student number of both members.
\begin{center}
   	\framebox{%
		\vbox{%
			\vskip 3mm
			\hbox to 6.28in{\rlap{\bf Name: Bharat Bhargava}\hfil{\bf Student Number: 1010380892}\hfil}
			\vskip 5mm
			\hbox to 6.28in{\rlap{\bf Name: Daivik Dhar}\hfil{\bf Student Number: 1010214260}\hfil}
			\vskip 3mm
		}%
	}
\end{center}
   
{\bf You should hand in:} 1) A scanned \textsf{.pdf} version of this sheet with your answers (file size
should be under $2$ MB); 2) Two figures for Part 1 Question 1 (c), four figures for Part 2 Question 2 and three figures for Part 2 Question 4
in the \textsf{.pdf} format; and 3) Two Python files \textsf{ldaqda.py} and \textsf{regression.py} that
contain your code. All these files should be uploaded to Quercus.

\section{Classification with Gaussian Models}
\begin{enumerate}
\item
\begin{enumerate}
\item Write down the maximum likelihood estimates of the parameters
$\boldsymbol \mu_{1}$, $\boldsymbol \mu_{2}$, $\boldsymbol \Sigma$,
$\boldsymbol \Sigma_1$, and $\boldsymbol \Sigma_2$ as functions of the training
data $\{\mathbf{x}_n, y_n\}, n=1,2,\ldots,N$. ({\bf $1$ pt})
\begin{center}
   \framebox{
    \begin{minipage}{5.5in}
   \vspace{4mm}
   
   \textbf{Class Means:} The ML estimates are simply the sample averages of the feature vectors for each respective class, filtered using the indicator function $I(\cdot)$.
   \begin{align*}
   \hat{\mu}_1 &= \frac{\sum_{n=1}^N I(y_n = 1) x_n}{\sum_{n=1}^N I(y_n = 1)}, \quad \hat{\mu}_2 = \frac{\sum_{n=1}^N I(y_n = 2) x_n}{\sum_{n=1}^N I(y_n = 2)}
   \end{align*}
   
   \textbf{QDA Covariances:} For QDA, each class has its own distinct covariance matrix, calculated using only the squared deviations from that specific class.
   \begin{align*}
   \hat{\Sigma}_1 &= \frac{\sum_{n=1}^N I(y_n = 1) (x_n - \hat{\mu}_1)(x_n - \hat{\mu}_1)^T}{\sum_{n=1}^N I(y_n = 1)} \\[6pt]
   \hat{\Sigma}_2 &= \frac{\sum_{n=1}^N I(y_n = 2) (x_n - \hat{\mu}_2)(x_n - \hat{\mu}_2)^T}{\sum_{n=1}^N I(y_n = 2)}
   \end{align*}
   
   \textbf{LDA Shared Covariance:} For LDA, both classes share the exact same covariance matrix. This is calculated as the pooled sample covariance averaged across all $N$ samples.
   \begin{align*}
   \hat{\Sigma} &= \frac{1}{N} \sum_{n=1}^N \Big[ I(y_n = 1) (x_n - \hat{\mu}_1)(x_n - \hat{\mu}_1)^T \\
   &\qquad\qquad\qquad + I(y_n = 2) (x_n - \hat{\mu}_2)(x_n - \hat{\mu}_2)^T \Big]
   \end{align*}
   
   \vspace{2mm}
   \end{minipage}
   }
   \end{center}
   \newpage

\item In the case of LDA, write down the decision boundary as a linear equation of
$\mathbf{x}$ with parameters $\boldsymbol \mu_{1}$, $\boldsymbol \mu_{2}$, and $\boldsymbol \Sigma$. Note that we assume $\pi=0.5$. ({\bf $0.5$ pt})
\begin{center}
   \framebox{
      \begin{minipage}{5.5in}
      \vspace{15mm}
      By equating the log-posteriors and assuming a shared covariance $\boldsymbol{\Sigma}$, the quadratic terms cancel, leaving the linear boundary:
      \begin{equation*}
      (\boldsymbol{\mu}_1 - \boldsymbol{\mu}_2)^T \boldsymbol{\Sigma}^{-1} \mathbf{x} - \frac{1}{2} \boldsymbol{\mu}_1^T \boldsymbol{\Sigma}^{-1} \boldsymbol{\mu}_1 + \frac{1}{2} \boldsymbol{\mu}_2^T \boldsymbol{\Sigma}^{-1} \boldsymbol{\mu}_2 = 0
      \end{equation*}
      \vspace{15mm}
      \end{minipage}
   }
\end{center}

In the case of QDA, write down the decision boundary as a quadratic equation of
$\mathbf{x}$ with parameters $\boldsymbol \mu_{1}$, $\boldsymbol \mu_{2}$, $\boldsymbol \Sigma_1$, and $\boldsymbol \Sigma_2$. Note that we assume $\pi=0.5$. ({\bf $0.5$ pt})
\begin{center}
   \framebox{
      \begin{minipage}{5.5in}
      \vspace{8mm}
      By equating the log-posteriors with distinct covariances $\boldsymbol{\Sigma}_1$ and $\boldsymbol{\Sigma}_2$, we obtain the quadratic boundary:
      \begin{align*}
      -\frac{1}{2} \mathbf{x}^T (\boldsymbol{\Sigma}_1^{-1} - \boldsymbol{\Sigma}_2^{-1}) \mathbf{x} &+ (\boldsymbol{\mu}_1^T \boldsymbol{\Sigma}_1^{-1} - \boldsymbol{\mu}_2^T \boldsymbol{\Sigma}_2^{-1}) \mathbf{x} \\
      &- \frac{1}{2} \boldsymbol{\mu}_1^T \boldsymbol{\Sigma}_1^{-1} \boldsymbol{\mu}_1 + \frac{1}{2} \boldsymbol{\mu}_2^T \boldsymbol{\Sigma}_2^{-1} \boldsymbol{\mu}_2 - \frac{1}{2}\ln \frac{|\boldsymbol{\Sigma}_1|}{|\boldsymbol{\Sigma}_2|} = 0
      \end{align*}
      \vspace{8mm}
      \end{minipage}
   }
\end{center}

\item Complete function \textsf{discrimAnalysis} in \textsf{ldaqda.py} to
	visualize LDA and QDA models and the corresponding decision boundaries. Please name the figures as \textsf{lda.pdf}, and \textsf{qda.pdf}. ({\bf $1$ pt})

\end{enumerate}
\item The misclassification rates are \framebox{\vbox{\vspace{3mm}\hbox to 16mm {\hfill 10.67\% \hfill}}} for LDA, and \framebox{\vbox{\vspace{3mm}\hbox to 16mm {\hfill 13.33\% \hfill}}} for QDA. ({\bf $1$ pt})
\end{enumerate}

\newpage

\section{Bayesian Linear Regression}
\begin{enumerate}
\item Express posterior distribution $p(\mathbf{a}|z_1,\ldots,z_{N} ; x_1, \ldots, x_N)$ using $\sigma^2, \beta$, $x_1, z_1, x_2, z_2, \ldots, x_N, z_N$.
   ({\bf $1$ pt})

   \begin{center}
\framebox{
   \begin{minipage}{5.5in}
   \vspace{6mm}
   
   Let $\mathbf{z} = [z_1, \dots, z_N]^T$ be the target vector, and $\mathbf{X}$ be the $N \times 2$ augmented input matrix where each row is $[1, x_n]$. By Bayes' rule, the posterior distribution is proportional to the likelihood times the prior:
   $$p(a | \mathbf{z}; \mathbf{x}) \propto p(\mathbf{z} | a; \mathbf{x}) p(a)$$
   
   Given the additive Gaussian noise $w \sim \mathcal{N}(0, \sigma^2 \mathbf{I})$, the likelihood is $p(\mathbf{z} | a; \mathbf{x}) = \mathcal{N}(\mathbf{X}a, \sigma^2 \mathbf{I})$. The prior is given as $p(a) = \mathcal{N}(0, \beta \mathbf{I})$. 
   
   Multiplying these two Gaussian distributions (by combining their exponents and completing the square) yields a new Multivariate Gaussian distribution for the posterior:
   $$p(a | z_1, \dots, z_N; x_1, \dots, x_N) = \mathcal{N}(\boldsymbol{\mu}_N, \boldsymbol{\Sigma}_N)$$
   
   Where the posterior covariance $\boldsymbol{\Sigma}_N$ and mean $\boldsymbol{\mu}_N$ are defined as:
   $$ \boldsymbol{\Sigma}_N = \left( \frac{1}{\beta} \mathbf{I} + \frac{1}{\sigma^2} \mathbf{X}^T \mathbf{X} \right)^{-1} $$
   $$ \boldsymbol{\mu}_N = \frac{1}{\sigma^2} \boldsymbol{\Sigma}_N \mathbf{X}^T \mathbf{z} $$
   
   \vspace{6mm}
   \end{minipage}
}
\end{center}
\item Let $\sigma^2=0.1$ and $\beta=1$. Draw four contour plots corresponding to the distributions $p(\mathbf{a})$, $p(\mathbf{a}|z_1;x_1)$, $p(\mathbf{a}|z_1,\ldots, z_{5}; x_1 \ldots x_5)$, and $p(\mathbf{a}| z_1,\ldots,z_{100}; x_1 \ldots x_{100})$. In all contour plots, the x-axis represents $a_0$, and the y-axis represents $a_1$. Please save the figures with names \textbf{prior.pdf, posterior1.pdf, posterior5.pdf, posterior100.pdf}, respectively. ({\bf $1.5$ pt})

\item Suppose that there is a new input $x$, for which we want to predict the corresponding target value $z$. Write down the distribution of the prediction $z$, i.e, $p(z|z_1,\ldots,z_N ; x, x_1, \ldots x_N)$. ({\bf $1$ pt})
  \begin{center}
\framebox{
   \begin{minipage}{5.5in}
   \vspace{6mm}
   Let the augmented new input be $\mathbf{x}_{new} = [1, x]^T$. The predicted target follows the linear model $z = \mathbf{x}_{new}^T a + w$, where the noise is $w \sim \mathcal{N}(0, \sigma^2)$.\\[6pt]
   Since the posterior distribution $p(a|\mathbf{z}; \mathbf{x}) = \mathcal{N}(\boldsymbol{\mu}_N, \boldsymbol{\Sigma}_N)$ and the noise $w$ are both independent Gaussians, their linear combination $z$ is also a Gaussian distribution. (by the property of Gaussian distributions) \\[6pt]
   The expected mean is $\mathbf{x}_{new}^T \boldsymbol{\mu}_N$ and the variance includes both the parameter uncertainty and the noise variance ($\mathbf{x}_{new}^T \boldsymbol{\Sigma}_N \mathbf{x}_{new} + \sigma^2$). Therefore, the predictive distribution is:
   \begin{equation*}
   p(z | z_1, \dots, z_N ; x, x_1, \dots x_N) = \mathcal{N} \left( \mathbf{x}_{new}^T \boldsymbol{\mu}_N, \ \mathbf{x}_{new}^T \boldsymbol{\Sigma}_N \mathbf{x}_{new} + \sigma^2 \right)
   \end{equation*}
   \vspace{2mm}
   \end{minipage}
}
\end{center}
 
\newpage

\item Let $\sigma^2=0.1$ and $\beta=1$. Given a set of new inputs $\{-4,-3.8,\ldots,3.8,4\}$, plot three figures, whose x-axis is the input and y-axis is the prediction, corresponding to three cases:
    \begin{enumerate}
    \item The predictions are based on one training sample, i.e., based on $p(z|z_1;x,x_1,)$.
    \item The predictions are based on $5$ training samples, i.e., based on $p(z|z_1,\ldots,z_{5};x,x_1,\ldots,x_{5})$.
    \item The predictions are based on $100$ training samples, i.e., based on  $p(z|z_1,\ldots,z_{100};x,x_1,\ldots,x_{100})$.
    \end{enumerate}

     The range of each figure is set as $[-4, 4]\times [-4, 4]$. Each figure should contain the following three components: 1) the new inputs and the corresponding predicted targets; 2) a vertical interval at each predicted target, indicating the range within one standard deviation; 3) the training sample(s) that are used for the prediction. Use \textsf{plt.errorbar} for 1) and 2); use \textsf{plt.scatter} for 3). Please save the figures with names \textbf{predict1.pdf, predict5.pdf, predict100.pdf}, respectively. ({\bf $1.5$ pt})

\end{enumerate}

\end{document}
